\documentclass{aitthesis}
\usepackage[utf8]{inputenc}
\usepackage[T1]{fontenc}
\usepackage{geometry}
\usepackage{graphicx}
\usepackage{float}
\usepackage{hyperref}
\usepackage[english,indonesian]{babel}

\title{Optimasi Pemrosesan Template LaTeX Menggunakan Pendekatan Berbasis AI}
\author{Academic Support Agent}
\date{\today}

\begin{document}
\maketitle

\begin{abstract}

\end{abstract}

\chapter{PENDAHULUAN}


\chapter{TINJAUAN PUSTAKA}


% Original content (commented out):
% \documentclass{aitthesis}
\usepackage{amsmath,amssymb}

\usepackage{multirow}
\usepackage{graphicx}
\usepackage{subfig}

\bibliographystyle{apacite}

\begin{document}

\pagenumbering{roman}
\addtocontents{toc}{~\hfill\textbf{Page}\par} %This adds the word Page on top of the table of content

\include{cover}
\include{declaration}
\include{acknowledgements}
\include{abstract}
\include{contents}
\include{loft}

\pagestyle{plain}
\pagenumbering{arabic}

\include{introduction}
\include{literature}
\include{methodology}
\include{results}
\include{conclusion}

\newpage
\phantomsection
\addcontentsline{toc}{chapter}{\bf REFERENCES}
\bibliography{references}

\newpage\pagestyle{plain}

\include{appendix}

\begin{table}
\caption{An example table in latex.}
\begin{center}
\begin{tabular}{l l}
\hline
    Methods & Metric\\ \hline
Method A      & 153.3                \\ 
Method B & 2.4                  \\ \hline
\end{tabular}
\label{tab:monkey}
\end{center}
\small{\textit{Note.} Add notes here.}
\end{table}

\end{document}


\end{document}